\documentclass[12pt,a4paper]{article}
\usepackage[utf8]{inputenc}
\usepackage[spanish]{babel}
\usepackage{geometry}
\usepackage{listings}
\usepackage{xcolor}
\usepackage{parskip}
\usepackage{amsmath}
\usepackage{hyperref}

\geometry{margin=2.5cm}

\lstset{
  basicstyle=\ttfamily\small,
  keywordstyle=\color{blue}\bfseries,
  commentstyle=\color{gray},
  breaklines=true,
  frame=single,
  language=Java,
  numbers=left,
  numberstyle=\tiny\color{gray},
  literate={á}{{\'a}}1 {é}{{\'e}}1 {í}{{\'\i}}1 {ó}{{\'o}}1 {ú}{{\'u}}1 {ñ}{{\~n}}1 {Á}{{\'A}}1 {É}{{\'E}}1 {Ó}{{\'O}}1
}

\title{\textbf{Sopa de Letras con Resolución Automática}\\
  \large Documentación Técnica}
\author{Inteligencia Artificial}
\date{}

\begin{document}
\maketitle
\tableofcontents
\newpage

\section{Descripción del Proyecto}
Juego de \textbf{Sopa de Letras} desarrollado en Java con Swing y dark theme. Las palabras se cargan desde un archivo externo \texttt{palabras.txt}. Incluye resolución automática con animación, temporizador por palabra y generación dinámica de nuevas sopas.

\section{Generación de la Sopa}

\subsection{Cargar palabras}
Las palabras se leen del archivo \texttt{palabras.txt}, una por línea, y se almacenan en una lista. Se normalizan a mayúsculas y se eliminan tildes para simplificar la comparación.

\subsection{Colocar palabras en el tablero}
Para cada palabra se intentan posiciones y direcciones al azar hasta encontrar un espacio válido:

\begin{lstlisting}
boolean colocarPalabra(String palabra) {
    int intentos = 0;
    while (intentos++ < 1000) {
        int fila = rand.nextInt(FILAS);
        int col  = rand.nextInt(COLS);
        int dir  = rand.nextInt(8); // 8 direcciones

        if (cabe(palabra, fila, col, dir) &&
            noColisiona(palabra, fila, col, dir)) {
            escribir(palabra, fila, col, dir);
            return true;
        }
    }
    return false; // no pudo colocarse
}
\end{lstlisting}

Las 8 direcciones son: derecha, izquierda, abajo, arriba y las 4 diagonales.

\subsection{Rellenar espacios vacíos}
Después de colocar todas las palabras, las celdas vacías se llenan con letras aleatorias del abecedario español para ocultar las palabras.

\section{Algoritmo de Búsqueda (Solver)}

El solver busca cada palabra verificando todas las celdas del tablero como posible punto de inicio:

\begin{lstlisting}
int[] buscarPalabra(String palabra) {
    int[][] dirs = {{0,1},{0,-1},{1,0},{-1,0},
                   {1,1},{1,-1},{-1,1},{-1,-1}};
    for (int f = 0; f < FILAS; f++) {
        for (int c = 0; c < COLS; c++) {
            for (int[] dir : dirs) {
                if (coincide(palabra, f, c, dir))
                    return new int[]{f, c,
                                     dir[0], dir[1]};
            }
        }
    }
    return null;
}
\end{lstlisting}

La función \texttt{coincide()} compara carácter por carácter en la dirección indicada, verificando que no salga del tablero.

\section{Resaltado y Animación}

Cuando se encuentra una palabra (manual o automáticamente):
\begin{itemize}
  \item Se calculan todas las celdas que ocupa.
  \item Se pintan en color destacado usando \texttt{paintComponent()}.
  \item La palabra se marca como encontrada en la lista lateral.
  \item Para \textbf{Resolver Todo}, se usa un \texttt{javax.swing.Timer} con 300ms entre palabras para crear un efecto de animación.
\end{itemize}

\section{Detección Manual}

El jugador hace clic en la primera letra de la palabra y arrastra hasta la última. Se registran las coordenadas de inicio y fin del arrastre, y se valida si forman una línea recta en alguna de las 8 direcciones. Si la palabra formada coincide con alguna de la lista, se resalta automáticamente.

\section{Temporizador}

Cada palabra tiene un tiempo límite configurable. Un \texttt{javax.swing.Timer} decrementa el contador cada segundo y actualiza la etiqueta de tiempo. Si el tiempo llega a 0, se pasa automáticamente a la siguiente palabra.

\section{Interfaz Gráfica}

\begin{itemize}
  \item \textbf{Panel de tablero}: dibujado con \texttt{paintComponent()}, fondo oscuro (\texttt{\#1e1e1e}), letras en blanco.
  \item \textbf{Panel lateral}: lista de palabras a encontrar con estado visual (tachadas las encontradas).
  \item \textbf{Botones}: ``Nueva Sopa'', ``Siguiente Palabra'', ``Resolver Todo''.
  \item Mouse listeners detectan inicio y fin del arrastre del jugador.
\end{itemize}

\end{document}
